% Generated by src/targets.py and updated for v4 revision (referee comments 4, 16).
% Nominal-proxy table. Intervals are bounded ranges over the stated albedo
% and radius-factor assumptions, NOT confidence levels; rows marked (a) have
% M sin i above the Chen & Kipping (2017) Terran/Neptunian transition.
\begin{table*}[t]
\centering
\caption{\rev{Nearest confirmed temperate planets (plus one solar-type RV
signal whose planet interpretation has since been refuted) as SGL imaging
targets -- explicitly a \emph{nominal-proxy} 
table, not a propagated uncertainty budget. Radii use the \citet{chenkipping17} Terran branch $R_p = 1.01\,(M\sin i/M_\oplus)^{0.279}\,R_\oplus$ evaluated at $M\sin i$ (no de-inclination correction is folded into the nominal value; conditional isotropic factors are given in the text); $\pm$ values are the $4.03\%$ intrinsic dispersion of that branch. Host mass $M_\star$ and semimajor axis $a$ are listed to enable verification of the Keplerian transverse velocity budget $\Delta v_{90}$ (eq.~\ref{eq:dvrule}) and insolation $S$. $\mathrm{SNR_C}$ follows the scalar rule of eq.~\ref{eq:snrrule} and the bracket is the bounded product of $A\in[0.15,0.45]$ with a $\pm15\%$ radius allowance -- an assumption, not a confidence interval. $\Delta v_{90}$ is the time-integral of projected transverse acceleration over 90\,d for a face-on orbit, with the corresponding edge-on value in parentheses. Observing geometry $z=650$\,AU, $d=1$\,m, $\lambda=1\,\mu$m, $t_s=1800$\,s.}}
\label{tab:targets}
\scriptsize
\setlength{\tabcolsep}{1.8pt}
{\bfseries
\resizebox{\textwidth}{!}{%
\begin{tabular}{lccccccccccccc}
\toprule
Planet & Discovery & Host (SpT) & $z_0$ & $M_\star$ & $a$ & $M\sin i$ & $R_p$\,$\pm$\,disp.$^{d}$ & $S$ & $P$ & $D_{\rm img}$ & $\mathrm{SNR_C}$\,[range]$^{b}$ & $\beta_{\rm ecl}$ & $\Delta v_{90}$ (face)$^{b}$ \\
& Reference & & (pc) & ($M_\odot$) & (au) & ($M_\oplus$) & ($R_\oplus$) & ($S_\oplus$) & (d) & (km) & (1800\,s) & (deg) & ($\mathrm{km\,s^{-1}}$, edge) \\
\midrule
Proxima Cen b & \citet{anglada16} & Proxima Centauri (M5.5V) & 1.30 & 0.122 & 0.0486 & 1.07 & 1.03\,$\pm$\,0.04 & 0.65 & 11.19 & 31.8 & 666\,[283--1148] & +44.8 & 5.81 (3.70) \\
Ross 128 b & \citet{bonfils18} & Ross 128 (M4V) & 3.38 & 0.168 & 0.0496 & 1.35 & 1.10\,$\pm$\,0.04 & 1.38 & 9.87 & 13.1 & 581\,[247--1003] & +0.5 & 2.94 (1.87) \\
GJ 1061 d & \citet{dreizler20} & GJ 1061 (M5.5V) & 3.67 & 0.120 & 0.0540 & 1.64 & 1.16\,$\pm$\,0.05 & 0.69 & 13.03 & 12.7 & 282\,[120--487] & +60.8 & 1.68 (1.07) \\
Teegarden c & \citet{dreizler24} & Teegarden's Star (M7V) & 3.83 & 0.089 & 0.0443 & 1.11 & 1.04\,$\pm$\,0.04 & 0.37 & 11.41 & 10.9 & 130\,[55--224] & -0.3 & 1.72 (1.10) \\
GJ 273 b\textsuperscript{a} & \citet{astudillo17} & GJ 273 (M3.5V) & 3.80 & 0.290 & 0.0911 & 2.89 & 1.36\,$\pm$\,0.06 & 1.06 & 18.65 & 14.4 & 490\,[208--846] & +16.5 & 1.34 (0.85) \\
$\tau$ Cet e\textsuperscript{a}\textsuperscript{e} & \citet{feng17} & $\tau$ Ceti (G8.5V) & 3.60 & 0.783 & 0.5380 & 3.93 & 1.48\,$\pm$\,0.06 & 1.71 & 162.90 & 16.5 & 909\,[386--1568] & +24.8 & 0.11 (0.07) \\
\bottomrule
\end{tabular}
}
}

\vspace{1mm}
\parbox{0.96\textwidth}{\footnotesize
\textbf{$^{a}$}\,$M\sin i$ lies above the \citet{chenkipping17} Terran/Neptunian transition ($2.04^{+0.66}_{-0.59}\,M_\oplus$), so the rocky-branch radius used here is an extrapolation, not a measurement.
\textbf{$^{b}$}The $\mathrm{SNR_C}$ bracket is a \emph{bounded range} over the stated albedo ($A\in[0.15,0.45]$) and radius allowance ($\pm15\%$), not a confidence level and not propagated from catalogue uncertainties. The $\Delta v_{90}$ entries give the face-on value with the corresponding edge-on value in parentheses; the spread over the starting mean anomaly inside a 90-day window is below $0.02\,\mathrm{km\,s^{-1}}$ for every row except $\tau$~Cet~e ($0.09$--$0.13\,\mathrm{km\,s^{-1}}$).
\textbf{$^{c}$}$\beta_{\rm ecl}$ is the ecliptic latitude of the antipodal focal direction, which sets the departure plane-change cost.
\textbf{$^{d}$}Radii, stellar parameters and orbital elements are from the per-row discovery catalogues referenced; none of these planets transits, so $M\sin i$ (not $M$) is the observed quantity.
\textbf{$^{e}$}\,$\tau$~Cet~e is listed for dynamics illustration only. The NASA Exoplanet Archive carries its 162.9-day signal as a \emph{false positive} after \citet{figueira25}, who attribute the $\tau$~Ceti RV signals to stellar activity; it is therefore not a confirmed target and its radius, $S$, SNR and $\Delta v$ inherit the assumed $M\sin i = 3.93\,M_\oplus$ of a refuted ephemeris.
}
\end{table*}
